\section{The defining-prime order and the exact coefficient specialization}
The received defining-prime account \cite[eqs.\ (9)--(31)]{DefiningPrimeReceived}
provides two integral ES witnesses and proposes the signed-order comparison.
The independent derivation below proves its arithmetic assertions with all
units retained. Here we give the full residue/conductor matrix and the exact
map from the coefficient extension ET1--39 to the actual witness. The
original finite inverse FS13 and residue RD1--4 retain their human
polynomial and residue provenance \cite{ES488,Tao,JouveRV}.

\subsection{All eight integral conductor generators}
Let $p\equiv1\pmod {12}$ be prime, and let positive integers $x,y,z$
satisfy $4/p=1/x+1/y+1/z$. Set
\[
 (r_0,r_1,r_2,r_3)=(p,x,y,z),\quad
 S=p+x+y+z,\quad A=-S^{-1},\quad
 f(r)=\prod_{i=0}^3(r-r_i),\quad h=Af .
 \tag{PS1}
\]
The four roots are distinct by the elementary ES argument in the
independent integral derivation. For now assume $p\nmid S$, so $A$
is a unit in $R_p=\mathbb Z_p$, with fraction field $K_p=\mathbb Q_p$.
Put
\[
 \begin{aligned}
 d_i&=Af'(r_i),&
 b_i&=v_p(d_i)=\sum_{j\ne i}v_p(r_i-r_j),\\
 m_i&=\lfloor b_i/2\rfloor,&
 \epsilon_i&=b_i-2m_i,&u_i&=p^{-b_i}d_i .
 \end{aligned}
 \tag{PS2}
\]
The symbols $d_i$ retain their full units; $d_i$ is not replaced by $p^{b_i}$.
The original integral algebra and its integral closure in its generic
fibre are
\[
 E=R_p[r,\sigma]/(f,\sigma^2-h'),\qquad
 O=\prod_{i=0}^3R_p[\tau_i],\qquad
 \tau_i^2=p^{\epsilon_i}u_i,\quad \sigma_i=p^{m_i}\tau_i .
 \tag{PS3}
\]
Each factor can be a quadratic field order or a split quadratic algebra.
The independent proof verifies its full integral closure directly, including
the split case. With $V=(r_i^j)_{0\le i,j\le3}$, the inclusion is
\[
 C_p=\diag(V,\diag(p^{m_i})V)
 \tag{PS4}
\]
in the unchanged basis $(1,r,r^2,r^3,\sigma,\sigma r,\sigma r^2,\sigma r^3)$
and the normalization basis $(1_0,\ldots,1_3,\tau_0,\ldots,\tau_3)$.
Thus every label and both geometric signs are retained.

Define $f_i=f/(r-r_i)$ and let $Q$ be the four-by-four matrix with
$i$th column the increasing-power coefficients of $f_i$:
\[
 Q_i=\begin{pmatrix}
 -\prod_{j\ne i}r_j\\
 \sum_{\substack{j<k\\j,k\ne i}}r_jr_k\\
 -\sum_{j\ne i}r_j\\1
 \end{pmatrix},\quad
 VQ=\diag(f'(r_i)),\quad\det Q=\det V .
 \tag{PS5}
\]
The evaluation identity proves the middle formula. Since
$\prod_i f'(r_i)=(-1)^6\prod_{i<j}(r_i-r_j)^2=(\det V)^2$,
taking determinants proves the last equality, with this precise quartic sign.

The exact algebra conductor $I=(E:O)=\{z\in O:zO\subset E\}$ is the
$R_p$-lattice with these eight original-coordinate generators:
\[
 \boxed{A p^{m_i}f_i(r),\qquad A\sigma f_i(r),\quad0\le i\le3.}
 \qquad
 J_I=\diag(AQ\diag(p^{m_i}),AQ).
 \tag{PS6}
\]
To prove equality with the conductor, retain the unchanged residue
\[
 \Lambda(R+\sigma T)=A^{-1}[r^3](T\bmod f),\quad
 \Lambda(R+\sigma T)=\sum_i\frac{T(r_i)}{d_i}.
 \tag{PS7}
\]
The second equality is Lagrange interpolation: the coefficient of $r^3$
in $f_i/f'(r_i)$ is $1/f'(r_i)$. The full Gram $B_E$ is exactly RD3,
whose determinant $A^{-8}$ is a unit. All of its entries and its inverse
are integral here. Hence $E$ is its own residue-dual lattice:
$E^\#=\{z:\Lambda(zE)\subset R_p\}=E$.

On the $i$th factor of $O$, the residue Gram in $(1_i,\tau_i)$ is
\[
 B_{O,i}=\begin{pmatrix}
 0&(p^{m_i}d_i)^{-1}\\(p^{m_i}d_i)^{-1}&0
 \end{pmatrix}.
 \tag{PS8}
\]
Indeed only the odd component contributes, and
$\tau_i=\sigma_i/p^{m_i}$. Thus
$O^\#=\prod_i(p^{m_i}d_i)O_i$.
For $z\in O^\#$ and $o\in O$, every $e\in E$ gives
$\Lambda(zoe)\in R_p$ because $oe\in O$; hence $zo\in E^\#=E$.
Conversely $zO\subset E$ implies $\Lambda(zO)\subset R_p$.
Therefore
\[
 I=O^\#=\prod_i(p^{m_i}d_i)O_i
       =\prod_i p^{\,3m_i+\epsilon_i}u_iO_i .
 \tag{PS9}
\]
The two generators on branch $i$ pull back by the labelled Lagrange
polynomial $L_i=f_i/f'(r_i)$ to
$p^{m_i}d_iL_i=Ap^{m_i}f_i$ and
$d_i\sigma L_i=A\sigma f_i$. This proves PS6 with every scalar unit.

More than the index is determined: the full Smith factors of $E/I$
are the same as those of $O/E$. Let $B_O$ be PS8 in the order PS4.
Then
\[
 B_E=C_p^{\mathsf T}B_OC_p,\qquad
 C_p^{-1}B_O^{-1}=B_E^{-1}C_p^{\mathsf T}.
 \tag{PS10}
\]
The left matrix generates $I$ in the original basis. Since $B_E^{-1}$
is an integral invertible matrix, its Smith factors are those of
$C_p^{\mathsf T}$, hence those of $C_p$.
Equivalently the exact pairing
\[
 E/I\times O/E\longrightarrow K_p/R_p,\qquad
 (\bar e,\bar o)\longmapsto\Lambda(eo)\bmod R_p
 \tag{PS11}
\]
is perfect. It is well defined by $I=O^\#$ and $E=E^\#$.
Its left kernel is $I$ and its right kernel is $E$ by these same
equalities; the equal elementary divisors from PS10 prove perfectness
on each finite cyclic factor.

In particular, for $g=\sum_{i<j}v_p(r_i-r_j)$,
\[
 \begin{aligned}
 \ell(O/E)=\ell(E/I)&=2g+\sum_i m_i
                   =3g-\tfrac12\sum_i\epsilon_i,\\
 \ell(O/I)&=2\ell(O/E),\\
 \det J_I&=A^8(\det V)^2p^{\sum_i m_i}.
 \end{aligned}\tag{PS12}
\]
The first equality follows also by taking determinants in PS4 and PS6.
The branch length in PS9 is $2(3m_i+\epsilon_i)$, confirming the second.

\subsection{The exact rank-twenty-four to rank-eight specialization}
Keep the coefficient names $C_0,D_0$ for the actual witness, so
\[
 C_0=-5Axyz,\qquad D_0=Apxyz,\qquad D_0=-pC_0/5 .
 \tag{PS13}
\]
The ring constructions ET1--20 work over
$\mathbb Z[1/10,A,A^{-1},B,C]$, because their displayed matrices
use only these inverted units and their proofs use monic division.
Specialize $(A,B,C)$ to the actual values $(A,B_0,C_0)\in R_p^3$.
The resulting two rank-twenty-four orders have exactly the original
ET bases and the inclusion $K(C_0)\otimes I_8$. Selecting the given
witness is the following pair of surjective algebra maps:
\[
 \begin{aligned}
 \Pi_D&=(I_8,\ D_0I_8,\ D_0^2I_8),&
 \Pi_P&=(I_8,\ pI_8,\ p^2I_8),\\
 \Pi_P(K(C_0)\otimes I_8)&=\Pi_D,&
 C_p\Pi_P(K(C_0)\otimes I_8)&=C_p\Pi_D .
 \end{aligned}\tag{PS14}
\]
The $P$ substitution satisfies its original monic cubic because the
witness satisfies the ES relation. The $D$ substitution satisfies ET1
by the same identity. They are algebra maps on the whole signed algebra:
the $r,\sigma$ images retain exactly $f$ and $h'$ at this coefficient
point. Entrywise substitution of $D_0=-pC_0/5$ proves both matrix
identities; the leading $I_8$ proves surjectivity.

The full kernels are free of rank sixteen, with no lost components:
\[
 \begin{aligned}
 \ker\Pi_D&=\{(-D_0u-D_0^2v,u,v):u,v\in R_p^8\},\\
 \ker\Pi_P&=\{(-pu-p^2v,u,v):u,v\in R_p^8\}.
 \end{aligned}\tag{PS15}
\]
Indeed these formulas solve the displayed equations without division.
Thus the rank-twenty-four coefficient comparison and the rank-eight
integral comparison have an exact commuting specialization, rather
than an identification of their different cokernels.

Its integral-ideal effect is fully calculable. Put $c=v_p(C_0)$.
The coefficient conductor $C_0^2 E_N$ specializes under $\Pi_P$ to
$C_0^2E$. Its inclusion in the arithmetic conductor $I$ holds exactly when
\[
 \boxed{2c\ge\max_i(3m_i+\epsilon_i).}
 \tag{PS16}
\]
Since $I$ is an $E$-ideal, the inclusion is equivalent to $C_0^2\in I$.
Branchwise PS9 turns that assertion into precisely the displayed
inequalities; the unit $u_i$ does not change a valuation.
In the one-divisible case $c=1$ and the other unit-root gap has
valuation $n=v_p(x-y)$; the branch exponents are
$1,1,3\lfloor n/2\rfloor+(n\bmod2)$ twice, up to the labelled order.
Hence PS16 holds exactly for $n\le1$.
In the two-divisible case $c=2$ and the three-root cluster has two gaps
of valuation one and one gap of valuation $n\ge1$.
Its exponents are zero at the isolated unit root, three at the separated
cluster root, and
$3\lfloor(n+1)/2\rfloor+((n+1)\bmod2)$ at both closer roots.
Hence PS16 holds exactly for $n\le2$.
These are complete alternatives on the classified ES strata.
Failure of this inclusion for larger gaps is an exact obstruction
between these two particular conductor ideals, while PS14 continues
to give their full algebraic relation.

\subsection{The two actual witnesses, all labels and signs}
At $p=1201$ the two input witnesses are
\[
 (x,y,z)=(306,16218,1082101),\qquad
 (x,y,z)=(306,21618,61251).
 \tag{PS17}
\]
For the first,
$1/16218+1/1082101=23/(306\cdot1201)$; adding $1/306$
gives $4/1201$. For the second, the last two denominators are
$1201\cdot18,1201\cdot51$ and $1/18+1/51=23/306$, giving
the same exact equality. The literal sums are $1099826$ and $84376$.
In the ordered roots $(p,x,y,z)$ their complete coefficient vectors are
\[
 \begin{aligned}
 u^{(1)}={}&\left(
 -\frac1{1099826},-\frac{19205048257}{1099826},
 \frac{13425378223770}{549913},
 -\frac{3224775849349554}{549913}\right),\\
 u^{(2)}={}&\left(
 -\frac1{84376},-\frac{1449375207}{84376},
 \frac{506477475135}{21094},
 -\frac{121655889527427}{21094}\right).
 \end{aligned}\tag{PS18}
\]
Multiplying the four factors in PS1 gives every entry.
Modulo $1201$ these are $(410,100,0,0)$ and $(522,0,0,0)$.
The complete valuation and branch-unit lists are
\[
 \begin{array}{c|c|c|c|c}
 & (b_i)&(m_i)&(\epsilon_i)&(u_i\bmod1201)\\ \hline
 1&(1,0,0,1)&(0,0,0,0)&(1,0,0,1)&(75,93,581,1126)\\
 2&(2,0,2,2)&(1,0,1,1)&(0,0,0,0)&(850,42,640,449).
 \end{array}\tag{PS19}
\]
Here $u_i$ in the last column is precisely $p^{-b_i}Af'(r_i)$,
not the coefficient-vector notation in PS18.
These lists can be checked by the exact product
$Af'(r_i)=A\prod_{j\ne i}(r_i-r_j)$; every difference is an explicit
integer from PS17. The Smith exponent lists of both $O/E$ and $E/I$ are
\[
 (0,0,0,0,0,0,1,1),\qquad
 (0,0,0,1,1,2,2,3),
 \tag{PS20}
\]
respectively. The branch conductor exponents are $(1,0,0,1)$ and
$(3,0,3,3)$. The entire original generator matrix in either example
is PS5--6 evaluated at the four stated integers; these formulas
specify all sixty-four rational entries without deleting any direction.
The accompanying exact certificate additionally records every entry.

At the common reduced root zero, the two full original primary algebras are
\[
 \mathbb F_{1201}[r,\sigma]/(r^2,\sigma^2-200r),\qquad
 \mathbb F_{1201}[r,\sigma]/(r^3,\sigma^2-3r^2).
 \tag{PS21}
\]
They follow from $h'$ reduced modulo $r^2$ and $r^3$,
respectively, retaining the coefficients $410,100$ and $522$.
The remaining simple roots and signs are still the other PS3 factors;
PS21 does not stand for the entire eight-dimensional algebra.

Choose $i\in K_p$ with $i^2=-1$, possible since $p\equiv1\pmod4$,
and for each labelled root choose both signs
$\sigma_{i,\pm}=\pm\sqrt{d_i}$ in its quadratic algebraic closure.
Every original affine inverse coordinate is retained by
\[
 \begin{aligned}
 q_{i,\pm}=\big(&\sigma_{i,\pm}^{-1},\
 -r_i-i\sigma_{i,\pm},\
 A\sigma_{i,\pm}^3+2r_i\sigma_{i,\pm}+3i\sigma_{i,\pm}^2,\\
 &7ir_i^2\sigma_{i,\pm}
 +(B_0-17r_i+Ar_i^2)\sigma_{i,\pm}^2
 -13i\sigma_{i,\pm}^3-2A\sigma_{i,\pm}^4\big).
 \end{aligned}\tag{PS22}
\]
FS13 verifies $P(q_{i,\pm})=(A,B_0,C_0,D_0)$ exactly.
All $d_i$ are nonzero in characteristic zero, so all eight points exist.
Their first-coordinate valuations are $-b_i/2$:
in the first example $(-1/2,0,0,-1/2)$ and in the second
$(-1,0,-1,-1)$, for each of the two signs.
The integral boundary therefore records escaping affine coordinates
without changing the original polynomial Jacobian $-2$.

\subsection{The same rational algebra in the two original complex receivers}
For these rational witnesses define $E_{\mathbb Q}$ by PS3 over
$\mathbb Q$ with its literal rational $A,B_0,C_0,D_0$.
Its two base changes to $K_p$ and $\mathbb C$ retain the very same
eight coefficient basis vectors. The first has the maps PS4--22.
For the second, use exactly $F_U,F_W$ of GD14.
Thus, for example, the full arithmetic conductor generator matrix has
the two original complex realizations
\[
 \mathcal I_U=F_UJ_I,\qquad
 \mathcal I_W=F_WJ_I,\qquad
 (T_A\oplus T_A)\mathcal I_U=\mathcal I_W .
 \tag{PS23}
\]
The columns are literally
$(T_A^{-1}(Ap^{m_i}f_i)(S'-\kappa),0)$ and
$(0,T_A^{-1}(Af_i)(S'-\kappa))$ at the source, with the
$T_A^{-1}$ removed at the target. Their exact squared Gram matrices are
\[
 J_I^*\Gamma_UJ_I,\qquad J_I^*\Gamma_WJ_I,\quad
 \det(J_I^*\Gamma_bJ_I)
   =|A|^{16}|\det V|^4p^{2\sum m_i}\det\Gamma_b .
 \tag{PS24}
\]
Here $\Gamma_b=\diag(G_{b,r},G_{b,r})$ contains the entire original
moments and Gamma masses, and the displayed determinant is PS12
times its complex conjugate. It is nonzero for both actual witnesses.
These are exact complex embeddings of an integral lattice, not an
assignment of a $p$-adic valuation to a transcendental Gamma mass.

The coefficient specialization itself also has a complete complex
minimum in those same metrics. For either receiver set
$\Gamma=\Gamma_b$ and
\[
 \Pi_D=(I,D_0I,D_0^2I),\qquad
 d_D=1+|D_0|^2+|D_0|^4 .
\]
For every prescribed column $v\in\mathbb C^8$, the exact original
quotient minimum and its unique orthogonal representative are
\[
 \min_{\Pi_D z=v}\sum_{a=0}^2 z_a^*\Gamma z_a
       =\frac{v^*\Gamma v}{d_D},\qquad
 z_{\min}=\frac1{d_D}(v,\bar D_0v,\bar D_0^{\,2}v).
 \tag{PS25}
\]
Multiplication proves the constraint. For a kernel vector
$w$ in PS15, the inner product with this $z_{\min}$ is
$d_D^{-1}v^*\Gamma(w_0+D_0w_1+D_0^2w_2)=0$.
Pythagoras proves the minimum, uniqueness and the exact excess
$\|z\|^2=\|z_{\min}\|^2+\|z-z_{\min}\|^2$.
The identical proof for $\Pi_P$ replaces $D_0$ by the actual prime $p$.
Thus PS14 and PS25 give both the algebraic specialization and the
attained full-kernel minimum, with the original metric retained.

For the source-to-target specialization
$F_W\Pi_D(I_3\otimes F_U^{-1})$, the eight nonzero singular values
are exactly $\sqrt{d_D}$ times the eight original singular values
of $T_A\oplus T_A$, and its kernel is sixteen-dimensional.
Indeed its orthonormal representation is the tensor product of
the row $(1,D_0,D_0^2)$ with the original conductor representation;
the row has the sole positive singular value $\sqrt{d_D}$.
This provides the complete map relating the rank-twenty-four
coefficient cost to the rank-eight actual arithmetic witness.
It does not transfer a coefficient-lattice pole into a vanishing
original conductor moment.

